\chapter{Bases conceptuales}
\label{bases-conceptuales}
\section{Conceptos de teoría de la medida}
\begin{defin}
Sea $\mathcal{W}$ un conjunto. Se dice que $\mathcal{B} \subseteq \mathcal{P(\mathcal{W})}$ es una $\sigma$-álgebra de $\mathcal{W}$ si satisface:
\begin{enumerate}
\item $\emptyset \in \mathcal{B}$
\item $A \in B \implies A^c \in \mathcal{B}$
\item $A_n \in \mathcal{B}$ para todo $n \in \mathbb{N} \implies \bigcup_{n \in \mathbb{N}} A_n \in \mathcal{B}$ 
\end{enumerate}

La tupla $(\mathcal{W, B})$ se llama espacio medible. A veces se dice $\mathcal{W}$ medible.
\end{defin}

\begin{defin}
Sean $(\mathcal{W, B_W})$ y $(\mathcal{Z, B_Z})$ espacios medibles. Sea $f: \mathcal{W} \rightarrow \mathcal{B}$ una función. Diremos que $f$ es $\mathcal{B_W-B_Z}$ medible o simplemente medible si para todo $B \in \mathcal{B_Z}$ la preimagen $f^{-1}(B) \in \mathcal{B_W}$.
\end{defin}

\begin{defin}
Sea $(\mathcal{W, B})$ un espacio medible. Una medida $\mu$ en $(\mathcal{W, B})$ es una función
\begin{equation*}
\mu: \mathcal{B} \rightarrow \mathbb{R} \cup \{\infty\}
\end{equation*}
tal que
\begin{enumerate}
\item $\mu(A) \in [0, \infty]$
\item $\mu(\emptyset) = 0$
\item $\sigma-$aditividad: para toda secuencia de conjuntos $A_n \in \mathcal{B}$ disjuntos dos a dos:
\begin{equation*}
\mu\left(\bigcup_{n \in \mathbb{N}}A_n\right) = \sum_{n\in\mathbb{N}} \mu(A_n)
\end{equation*}
\end{enumerate}
\end{defin}

\section{Conceptos de probabilidad}
A lo largo del trabajo, asumimos siempre que estamos en un espacio de probabilidad $(\Omega, \mathcal{F}, P)$, con $\Omega$ el espacio muestral, $\mathcal{F}$ su $\sigma-$álgebra y $P:\mathcal{F} \rightarrow \mathbb{R}$ la función de medida de probabilidad que satisface los tres axiomas de Kolmogorov.

\begin{defin}[Variable aleatoria]
Una variable aleatoria es una función $X:\Omega \rightarrow \mathbb{R}$ tal que es medible, es decir, que cada preimagen de $X$ sobre un elemento de la $\sigma-$álgebra de Borel es un elemento de $\mathcal{F}$.
\end{defin}

\begin{defin}[Probabilidad condicionada]
Sean $A$ y $B$ dos sucesos tales que $P(B) > 0$. Definimos la probabilidad de $A$ condicionada a $B$ como
\begin{equation*}
P(A \mid B) = \frac{P(A \cap B)}{P(B)}
\end{equation*}
\end{defin}
Normalmente abreviaremos $P(A \cap B)$ o $P(X_1 \cap X_2 \cap \ldots\cap X_n)$ como $P(A, B)$ o $P(X_1, X_2, \ldots, X_n)$. En consecuencia de esta regla, tenemos $P(X_1, X_2, \ldots, X_n) = P(X_n\mid X_1, \ldots, X_{n-1})\ldots P(X_2\mid X_1)P(X_1)$ 
\begin{defin}[Independencia condicional]
Decimos que las variables $X$ e $Y$ son condicionalmente independientes dado $Z$ (denotado $(X \perp Y) \mid Z$ si
\begin{equation*}
P(X, Y \mid Z) = P(X\mid Z)P(Y\mid Z)
\end{equation*}
\end{defin}

\begin{teor}[Ley de la probabilidad total]
Supongamos $I \subseteq \mathbb{N}$ un conjunto finito o numerable. Digamos que los conjuntos $B_i$ con $i \in I$ son tales que:
\begin{itemize}
    \item $B_i \in \mathcal{F}$ para todo $i \in I$.
    \item $i \neq j \rightarrow B_i \cap B_j = \emptyset$.
    \item $\bigcup_{i \in I} B_i = \Omega$.
    \item $P(B_i) > 0$ para todo $i \in I$.
\end{itemize}
Entonces, para cualquier suceso $A \in \mathcal{F}$:
\begin{equation*}
P(A) = \sum_{i \in I} P(A\mid B_i) P(B_i)
\end{equation*}
\end{teor}
\begin{proof}
Como $A \subseteq \Omega$,
\begin{equation*}
    A = A \cap \Omega = A \cap \left(\bigcup_{i \in I} B_i\right)
\end{equation*}
Por la propiedad distributiva de la unión y la intersección, tenemos que 
\begin{equation*}
A \cap \left(\bigcup_{i \in I} B_i\right) = \bigcup_{i \in I} (B_i \cap A)
\end{equation*}

Aplicando ahora el axioma de la $\sigma-$aditividad de $P$ (utilizamos que los conjuntos $B_i$ son una partición de $\Omega$), tenemos que 

\begin{equation*}
P(A) = P\left(\bigcup_{i \in I} (B_i \cap A)\right) = \sum_{i \in I} P(B_i \cap A)
\end{equation*}

Aplicando ahora la definición de probabilidad condicionada, llegamos al resultado deseado:

\begin{equation*}
P(A) = \sum_{i \in I} P(B_i \cap A) = \sum_{i \in I} P(A \mid B_i)P(B_i)
\end{equation*}

\end{proof}

\section{Teoría de grafos}
\begin{defin}[Grafo dirigido]
Decimos que el par $G = \langle V, E\rangle$ donde 
\begin{itemize}
    \item $V = \{X_1, X_2, \ldots, X_n \mid n \in \mathbb{N}\}$
    \item $E \subseteq V \times V$ es un subconjunto de los pares ordenados con elementos en $V$.
\end{itemize}
es un grafo dirigido con vértices $V$ y aristas $E$.
\end{defin}

\textbf{Notación:} para todo par $(X_i, X_j) \in E$, diremos $X_i \rightarrow X_j$

\begin{defin}[Parentesco]
Para todo vértice del grafo $X_i$, definimos
\begin{itemize}
\item Los \textbf{padres} de $X_i$: $\text{PA}(X_i)=\{X_j \in V \mid X_j \rightarrow X_i \}$
\item Los \textbf{hijos} de $X_i$: $\text{CH}(X_i)=\{X_j \in V \mid X_i \rightarrow X_j \}$
\end{itemize}
\label{def:parentesco}
\end{defin}

\begin{defin}[Camino]
Dado un grafo dirigido $G=(V,E)$, un \textbf{camino} (o sendero) entre dos nodos $X$ e $Y$ es una secuencia de vértices distintos $(V_0,V_1,\ldots,V_k)$ tal que:
\begin{enumerate}
\item $V_0=X$ y $V_k=Y$.
\item Para cada $i \in \{0,\ldots,k-1\}$, los vértices son adyacentes en $G$. Es decir, se cumple $(V_i,V_{i+1}) \in E$ o bien $(V_{i+1},V_i)\in E$.
\end{enumerate}
\label{def:camino}
\end{defin}

\begin{defin}[Camino dirigido]
Dado un grafo dirigido $G=(V,E)$, un \textbf{camino dirigido} entre dos nodos $X$ e $Y$ es una secuencia de vértices distintos $(V_0,V_1,\ldots,V_k)$ tal que:
\begin{enumerate}
\item $V_0=X$ y $V_k=Y$.
\item Para cada $i \in \{0,\ldots,k-1\}$, los vértices están conectados en $G$. Es decir, se cumple $(V_i,V_{i+1}) \in E$.
\end{enumerate}
\end{defin}

\begin{defin}[Ciclo Dirigido]
Un \textbf{ciclo dirigido} es una secuencia de vértices $(V_0, V_1, \ldots, V_k)$ con $k \geq 1$ tal que:
\begin{enumerate}
    \item $(V_0, \ldots, V_{k-1})$ es un camino dirigido (vértices distintos).
    \item $V_0 = V_k$ (el vértice inicial coincide con el final).
\end{enumerate}
\end{defin}

\begin{defin}[Grafo Acíclico Dirigido - DAG]
Un grafo dirigido $G = (V, E)$ se dice \textbf{acíclico} si no contiene ningún ciclo dirigido.
A este tipo de estructuras se las conoce universalmente por sus siglas en inglés: \textbf{DAG} (\textit{Directed Acyclic Graph}).
\end{defin}